%Version 3.1 December 2024
\documentclass[pdflatex,sn-basic]{sn-jnl}

\usepackage{graphicx}
\usepackage{float}
\usepackage{multirow}
\usepackage{amsmath,amssymb,amsfonts}
\usepackage{amsthm}
\usepackage{mathrsfs}
\usepackage[title]{appendix}
\usepackage{xcolor}
\usepackage{textcomp}
\usepackage{manyfoot}
\usepackage{booktabs}
\usepackage{algorithm}
\usepackage{algorithmicx}
\usepackage{algpseudocode}
\usepackage{listings}
\usepackage{tabularx}

\theoremstyle{thmstyleone}
\newtheorem{theorem}{Theorem}
\newtheorem{proposition}[theorem]{Proposition}

\theoremstyle{thmstyletwo}
\newtheorem{example}{Example}
\newtheorem{remark}{Remark}

\theoremstyle{thmstylethree}
\newtheorem{definition}{Definition}

\raggedbottom

\begin{document}

\title[Experimental Comparative Analysis of Conventional Feature Extraction vs. Deep Learning Models for Facial Emotion Recognition]{Experimental Comparative Analysis of Conventional Feature Extraction vs. Deep Learning Models for Facial Emotion Recognition}

\author*[1]{\fnm{Hind} \sur{Laziri}}\email{laziri.h825@ucd.ac.ma}
\author[1]{\fnm{Mohammed Essaid} \sur{Riffi}}\email{saidriffi2@gmail.com}

\affil*[1]{\orgdiv{Department of Computer Science}, \orgname{Chouaib Doukali Faculty}, \orgaddress{\city{El Jadida}, \country{Morocco}}}

\abstract{Facial Emotion Recognition (FER) is a significant research area in human-computer interaction, with applications in healthcare and security. Historically, FER has been addressed using conventional methods based on handcrafted feature extraction, such as Local Binary Patterns (LBP) combined with Support Vector Machines (SVM). This study presents an experimental comparative analysis of three paradigms: a conventional LBP+SVM approach, a baseline CNN, and an optimized architecture, EfficientNet-B0, on the FER-2013 dataset. Experimental results indicate that while the LBP+SVM baseline achieves 27.96\% accuracy, the basic CNN reaches 62.16\%. The EfficientNet-B0 architecture demonstrates superior performance with an accuracy of 72.45\%, highlighting the effectiveness of optimized deep learning models. These approaches are evaluated in terms of accuracy, weighted F1-score, and computational efficiency, providing empirical benchmarks for selecting FER methods based on application constraints.}

\keywords{Facial Emotion Recognition (FER), Deep Learning, CNN, Feature Extraction, Comparative Study, Affective Computing}

\maketitle

\section{Introduction}\label{sec1}

\subsection{Context and Significance}
Emotions influence human thoughts, actions, and relationships. The ability to automatically recognize and interpret these expressions is the foundation of Facial Emotion Recognition (FER). FER is a component in human-computer interaction, affective computing, and real-world applications including healthcare, education, and security systems \cite{ref1}.

\subsection{Problem Statement and State of the Art}
FER remains a challenging task due to variations in head pose, illumination, occlusion, and the subtlety of human expressions. Historically, conventional methods relied on hand-crafted feature extraction techniques such as Local Binary Patterns (LBP) and Histograms of Oriented Gradients (HOG), followed by classifiers like Support Vector Machines (SVM) or K-Nearest Neighbors (K-NN) \cite{ref2}. While these methods are computationally efficient, their performance often decreases in unconstrained environments.

In recent years, deep learning architectures, particularly Convolutional Neural Networks (CNNs) and Vision Transformers (ViTs), have been applied to FER by learning hierarchical feature representations directly from raw images \cite{ref3, ref4}. These models achieve high accuracy but require annotated datasets and computational resources.

\subsection{Research Gap and Contribution}
The literature lacks a structured, comparative analysis of conventional versus deep learning approaches that quantifies the trade-offs between accuracy, robustness, and computational cost across diverse datasets. This paper addresses this gap by comparing these paradigms—ranging from handcrafted feature extraction to optimized deep architectures—and providing empirical evidence. By integrating EfficientNet-B0 alongside baseline models, this work quantifies the trade-offs between accuracy and computational cost, bridging the gap between theoretical performance and practical deployment constraints.

The source code and implementation details for this study are publicly available at \url{https://github.com/hindlaziri/FER-Review-Tradeoffs}.

The main contributions of this paper are:
\begin{itemize}
    \item A structured review methodology comparing conventional feature-based methods and deep learning architectures for FER.
    \item An analysis of widely used FER datasets, highlighting their inherent biases and their impact on model generalization.
    \item A quantitative trade-off analysis of accuracy versus computational cost between conventional and deep learning approaches, supported by a controlled baseline experimental validation on the FER-2013 dataset.
    \item Practical guidance for researchers and engineers in selecting FER methods tailored to specific application constraints.
\end{itemize}

\section{Methodology}\label{sec2}

\subsection{Methodology of Literature Review}
This review was inspired by PRISMA (Preferred Reporting Items for Systematic Reviews and Meta-Analyses) principles to ensure transparency. Literature was selected from academic databases including IEEE Xplore, ScienceDirect, SpringerLink, and ACM Digital Library. The search strategy utilized targeted keywords such as ''Facial Emotion Recognition'', ''Deep Learning FER'', ''CNN'', and ''FER Datasets''. The inclusion criteria focused on peer-reviewed journal articles published between 2019 and 2025. After applying the selection criteria, approximately 45 relevant studies were retained for analysis.

\subsection{Definition of Emotion and Facial Expression}
An emotion is defined as a complex reaction, encompassing both psychological and physiological dimensions. A facial emotion refers to the expression displayed on a person's face in response to an emotional state, resulting from facial muscle activity. The literature typically identifies seven basic emotions: Happiness, Sadness, Anger, Fear, Surprise, Disgust, and Neutral \cite{ref1}.

\subsection{Methodology: Preprocessing and Classification Techniques}

\subsubsection{Image Preprocessing Pipeline}
To ensure reliable model performance, raw facial images undergo a preprocessing pipeline aimed at standardizing inputs and maximizing feature quality. Figure 1 summarizes the typical end-to-end pipeline for Facial Emotion Recognition.

\begin{figure}[H]
    \centering
    \includegraphics[width=0.8\linewidth]{fer_pipeline.png}
    \caption{End-to-end pipeline for facial emotion recognition, illustrating the sequence of preprocessing steps from raw image input to emotion classification.}
    \label{fig:pipeline}
\end{figure}

The preprocessing stage typically involves face detection, alignment, normalization, and resizing to a fixed resolution (e.g., 48x48 or 224x224) required by the specific model architecture.

\subsubsection{Review of Recognition Methods}
Approaches to FER are generally categorized into conventional methods and deep learning architectures.

\begin{itemize}
    \item Conventional Methods: These methods rely on extracting pre-defined, hand-crafted features. Key techniques include Local Binary Patterns (LBP) and Histograms of Oriented Gradients (HOG) \cite{ref2}. These features are then fed into machine learning classifiers such as Support Vector Machines (SVM) \cite{ref3}.
    \item Deep Learning Approaches: These models learn hierarchical features directly from raw pixel data. A typical CNN architecture used for FER is illustrated in Figure 2.
\end{itemize}

\begin{figure}[H]
    \centering
    \includegraphics[width=0.8\linewidth]{cnn_architecture.png}
    \caption{Simplified Convolutional Neural Network (CNN) architecture for facial emotion recognition, demonstrating the progression through convolutional, pooling, and fully connected layers.}
    \label{fig:cnn_arch}
\end{figure}

\begin{itemize}
    \item Convolutional Neural Networks (CNNs): The standard for image analysis, capable of extracting complex hierarchical features (e.g., VGG, ResNet) \cite{ref5, ref16, ref17}. Recent surveys have highlighted the evolution of these architectures toward more efficient designs \cite{ref18, ref19}.
    \item Vision Transformers (ViTs): Advanced models utilizing self-attention mechanisms to capture global context, showing state-of-the-art performance but requiring massive data \cite{ref11, ref_vit}.
\end{itemize}

\section{Experimental Setup}\label{sec3}

\subsection{FER Datasets Overview}
The performance of FER models is dependent on the training data. Table 1 summarizes the characteristics of commonly used FER datasets.

\begin{table}[htbp]
\centering
\caption{Commonly Used FER Datasets}\label{tab:fer_datasets}
\begin{tabularx}{\linewidth}{@{}p{0.15\linewidth} p{0.15\linewidth} X X@{}}
\toprule
\textbf{Dataset} & \textbf{Size} & \textbf{Environment} & \textbf{Characteristics}\\
\midrule
CK+ \cite{ref8} & ~593 seq. & Controlled & High quality, posed, lab settings.\\
JAFFE \cite{ref10} & 213 images & Controlled & Japanese female faces, limited scale.\\
FER-2013 \cite{ref9} & ~35K images & In-the-wild & Auto-collected, label noise.\\
RAF-DB \cite{ref7} & ~30K images & In-the-wild & Diverse demographics, challenging pose/lighting.\\
AffectNet \cite{ref6} & >1M images & In-the-wild & Large scale, diversified labels.\\
\bottomrule
\end{tabularx}
\end{table}

\subsection{Quantitative Comparison: Accuracy and Computational Cost}

Table 2 presents a quantitative comparison of different approaches based on recent literature.

\begin{table}[htbp]
\centering
\caption{Quantitative Comparison of FER Approaches on Standard Datasets (Indicative Accuracy \%)}\label{tab:quantitative_comparison}
\begin{tabularx}{\linewidth}{@{}p{0.2\linewidth} c c c X X@{}}
\toprule
Approach & Acc. FER-2013 & Acc. CK+ & Cost & Advantages & Limitations\\
\midrule
LBP + SVM \cite{ref2} & 55 & 88 & Low & Simple & Low generalization.\\
HOG + K-NN \cite{ref3} & 60 & 90 & Low & Shape capture & Sensitive to light/pose.\\
CNN (e.g., ResNet) \cite{ref5} & 73 & 97 & High & Self-learning, robust. & Large datasets, memory.\\
Vision Transformer \cite{ref_vit} & 75 & 98 & Substantial & Global context & Intensive, overfitting.\\
\bottomrule
\end{tabularx}
\textit{Note: Direct comparison is indicative due to differing evaluation protocols across the cited studies.}
\end{table}

\subsection{Reproducibility Statement}
Experiments were conducted using PyTorch 1.12 with CUDA 11.6 on an NVIDIA Tesla V100 GPU, while conventional models were implemented using scikit-learn 1.0.

\subsection{Discussion and Error Analysis}

Conventional methods are utilized for their implementation simplicity and low computational cost. However, their reliance on handcrafted features limits their ability to capture complex emotional patterns. Deep learning models, particularly CNNs and ViTs, have established themselves as the choice for high-accuracy tasks. However, error analysis reveals common misclassifications, particularly between visually similar emotions such as Fear and Surprise \cite{ref15}. Furthermore, models trained on specific datasets often suffer from demographic bias \cite{ref14, ref20}. Addressing these biases remains a research area \cite{ref21, ref22}. Recent studies have also explored the use of synthetic data and data augmentation to mitigate these issues \cite{ref23, ref24, ref25}.

\subsection{Experimental Validation and Discussion}

To provide empirical evidence, we conducted a controlled baseline comparison evaluating a conventional approach (LBP + SVM) against two deep learning paradigms: a baseline CNN and an optimized architecture, EfficientNet-B0, on the FER-2013 dataset \cite{ref9}. EfficientNet-B0 \cite{tan2019efficientnet} was selected due to its compound scaling method, which offers a trade-off between accuracy and computational efficiency.

\subsection{Experimental Results and Comparative Analysis}

Table 3 summarizes the performance of the evaluated models on the FER-2013 test set.

\begin{table}[htbp]
\centering
\caption{Experimental Results on FER-2013 Test Set}\label{tab:results}
\begin{tabularx}{\linewidth}{@{}l c c c c@{}}
\toprule
Model & Accuracy (\%) & Weighted F1 (\%) & Train Time (s) & Inference Time (s)\\
\midrule
LBP + SVM & 27.96\% & 21.48\% & 160.9 & 28.76\\
Basic CNN & 62.16\% & 61.34\% & 1517.8 & 1.61\\
EfficientNet-B0 & 72.45\% & 71.92\% & 3420.5 & 2.45\\
\bottomrule
\end{tabularx}
\end{table}

The integration of EfficientNet-B0 yielded a performance leap, achieving a peak accuracy of 72.45\% and a weighted F1-score of 71.92\%. While the training duration for EfficientNet-B0 was higher (3420.5s), its inference time on the test set was 2.45s. This indicates that optimized architectures can provide high discriminative capability while maintaining the low-latency performance required for real-time applications.

The confusion matrices presented in Figure 4 demonstrate that the CNN model exhibits superior capability in discriminating between complex emotion categories. This advantage stems from the CNN's hierarchical feature learning. In contrast, the conventional LBP+SVM approach demonstrates lower performance, frequently misclassifying diverse emotions as 'Happy' or 'Neutral', because hand-crafted features lack the semantic depth necessary to differentiate subtle microexpressions in unconstrained environments.

Figure 5 presents the per-class F1-scores, demonstrating that the CNN outperforms the conventional method across all emotion categories. The conventional method demonstrated zero recognition capability for 'Disgust' (F1-score of 0.0\%), highlighting the limitations of hand-crafted features in capturing subtle facial micro-expressions. Furthermore, the overall performance of both models is constrained by inherent dataset biases in FER-2013, specifically severe class imbalance (e.g., the underrepresentation of the 'Disgust' class) and pervasive label noise resulting from automated collection methods. These generalization issues underscore the necessity for robust, diverse datasets and advanced techniques like domain adaptation to achieve reliable performance in unconstrained environments.

The Basic CNN achieved an accuracy of 62.16\%, outperforming the LBP+SVM approach (27.96\%). The conventional LBP+SVM approach demonstrates lower performance because hand-crafted features lack the semantic depth necessary to differentiate subtle microexpressions in unconstrained environments.

\section{Conclusion}\label{sec4}
This paper presented a comparative analysis of conventional, baseline deep learning, and optimized architectures for Facial Emotion Recognition. The findings suggest that while conventional methods remain viable for resource-constrained environments, deep learning approaches, particularly the EfficientNet-B0 architecture, offer superior performance for unconstrained scenarios. The inclusion of an optimized model in this study indicates that advanced architectures can mitigate the limitations of shallow networks and handcrafted features.

However, this study acknowledges certain limitations. The experimental validation focuses on the FER-2013 dataset, which introduces potential biases and prevents strict subject-independent evaluation. Consequently, models may learn demographic features rather than emotion-invariant patterns, which can limit generalization. Future research should prioritize cross-dataset validation and the mitigation of demographic biases to ensure robust deployment.

Furthermore, while Vision Transformers (ViTs) achieve high performance, their computational requirements were beyond the scope of this baseline study. Future work will focus on evaluating ViTs and hybrid CNN-Transformer models across multiple datasets, such as CK+, RAF-DB, and AffectNet, to establish broader performance benchmarks and further validate the generalization capability of deep learning approaches.

\nocite{*}
\bibliography{sn-bibliography}

\end{document}
